\documentclass[11pt,a4paper]{article}

\usepackage[utf8]{inputenc}
\usepackage[T1]{fontenc}
\usepackage[russian]{babel}
\usepackage{amsmath,amssymb,amsthm}
\usepackage{geometry}
\usepackage{hyperref}
\usepackage{mathrsfs}
\usepackage{parskip}

\geometry{left=2.5cm,right=2.5cm,top=3cm,bottom=3cm}

\title{
  \textbf{GRA-Chiral-Nullification-Math:} \\
  \textbf{Хиральность как нативный язык GRA-обнулёнки}
}

\author{
  oleg bits (\texttt{qqewq}) \\
  \texttt{https://github.com/qqewq} \\
  \texttt{GRA-Chiral-Nullification-Math} \\
  \texttt{https://github.com/qqewq/GRA-Chiral-Nullification-Math}
}

\date{}

\begin{document}

\maketitle

\begin{abstract}
В работе предлагается строгая математическая формализация GRA-обнулёнки в языке хиральности.
Хиральность трактуется не как свойство химических молекул, а как абстрактное свойство объектов в метрическом/топологическом пространстве (без химии).
Оператор обнулёнки \(\mathcal{N}\) интерпретируется как «хиральный подъём» политики в пространство пар (правое/левое) с последующим занулением весов и состояния.
Показывается, что GRA-обнулёнка и хиральность эквивалентны на уровне операторов: \( \mathcal{N} \equiv \text{хиральный подъём} \).
Работа предназначена для фундамента GRA-экосистемы: \( \text{GRA-Meta-Nulling-Foundations} \), \( \text{GRA-SubjectSwap} \), \( \text{GRA-ASI-Metric-Space} \) и др.
\end{abstract}

\section{Введение}

\subsection{GRA-обнулёнка и хиральность}

GRA-обнулёнка — не просто «поставить всё в ноль», а конструктивный переход

\[
(\Pi, W, S) \;\longrightarrow\; (\Pi', 0, 0),
\]

где:

- \(\Pi\) — политика (policy),
- \(W\) — вектор весов,
- \(S\) — состояние,
- \(\Pi' = F(\Pi)\) — переparameterизованная политика.

В этой работе мы показываем, что этот процесс можно **напрямую** трактовать как переход в **хиральное пространство**:

\[
(\Pi, W, S) \;\longrightarrow\; ((R, L), 0, 0),
\]

где:

- \(R = \Pi\) — «правая» конфигурация,
- \(L = Z(\Pi)\) — «левая» конфигурация, получаемая зеркальным оператором \(Z\),
- пара \((R, L)\) — хиральная пара.

Хиральность здесь не химия, а **абстрактное свойство объекта** не совпадать со своим зеркальным отражением через любые автоморфизмы.

\section{Математика хиральности}

\subsection{Пространство и зеркальный оператор}

\begin{definition}[Хиральность]
Пусть \(X\) — множество объектов, \(Z : X \to X\) — оператор зеркального отображения,
а \(\mathcal{R}\) — семейство допустимых автоморфизмов (например, вращений).

Элемент \(x \in X\) называется **хиральным**, если

\[
\forall \rho \in \mathcal{R} \quad \rho(x) \neq Z(x).
\tag{1}
\]
\end{definition}

\begin{definition}[Правая и левая конфигурации]
Для любого \(x \in X\) определяем:

\[
\mathcal{R}(x) = x, \qquad \mathcal{L}(x) = Z(x).
\tag{2}
\]

Пара \((\mathcal{R}(x), \mathcal{L}(x))\) называется **хиральной парой**.
\end{definition}

\begin{lemma}[Инволюция зеркала]
Если \(Z\) является инволюцией, то

\[
Z(Z(x)) = x \quad \text{для всех } x \in X.
\tag{3}
\]
\end{lemma}

\section{Математика GRA-обнулёнки}

\subsection{Тройка (политика, веса, состояние)}

\begin{definition}[Тройка GRA]
Пусть:

- \(\Pi \in \mathcal{P}\) — элемент пространства политик,
- \(W \in \mathcal{W}\) — вектор весов,
- \(S \in \mathcal{S}\) — состояние.

Тройка \((\Pi, W, S)\) называется **тройкой GRA**.
\end{definition}

\subsection{Оператор обнулёнки}

\begin{definition}[Оператор обнулёнки]
Оператор \(\mathcal{N}\) определяется как

\[
\mathcal{N} : \mathcal{P} \times \mathcal{W} \times \mathcal{S}
\to
\mathcal{P} \times \mathcal{W} \times \mathcal{S},
\]

\[
\mathcal{N}(\Pi, W, S) = (\Pi', 0, 0),
\tag{4}
\]

где \(\Pi' = F(\Pi)\) — переparameterизованная политика (reparametrization).
\end{definition}

\section{Эквивалентность: хиральность = обнулёнка}

\subsection{Хиральный подъём политики}

\begin{definition}[Хиральный подъём]
Для любой политики \(\Pi \in \mathcal{P}\) определяем **хиральный подъём** \(F\) как

\[
F(\Pi) = (\Pi, Z(\Pi)) = (R, L).
\tag{5}
\]
\end{definition}

\begin{theorem}[Обнулёнка как хиральный подъём]
Оператор обнулёнки \(\mathcal{N}\) можно переписать как

\[
\mathcal{N}(\Pi, W, S) = ((R, L), 0, 0),
\tag{6}
\]

где \(R = \Pi\), \(L = Z(\Pi)\), то есть

\[
\mathcal{N}(\Pi, W, S) = (F(\Pi), 0, 0).
\tag{7}
\]
\end{theorem}

\textbf{Доказательство.}
Из определений (4) и (5):

\[
\mathcal{N}(\Pi, W, S) = (\Pi', 0, 0)
= (F(\Pi), 0, 0)
= ((R, L), 0, 0).
\]

\(\square\)

\begin{corollary}[Обнулёнка = хиральность]
На уровне операторов GRA-обнулёнка и хиральность эквивалентны:

\[
\mathcal{N} \equiv \text{хиральный подъём}.
\tag{8}
\]
\end{corollary}

\section{Практические следствия}

\subsection{GRA-обнулёнка как хиральный слой}

В GRA-экосистеме обнулёнка может быть использована как:

- **хиральный слой** над субъект/инструмент (subject/instrument),
- способ формализовать «двойные» роли агента,
- инструмент для «обнулённого» переparameterизования политик в мета-обнуляющих системах.

См. также:

- \texttt{GRA-Meta-Nulling-Foundations},
- \texttt{GRA-SubjectSwap},
- \texttt{GRA-ASI-Metric-Space},
- \texttt{GRA-Quantum-Poetry}.

\section{Заключение}

Мы показали, что **GRA-обнулёнка** может быть строго описана как **хиральный подъём** в пространстве политик с последующим занулением весов и состояния.
Хиральность формализована чисто математически: без химии, без LLM, только объекты, зеркальный оператор и автоморфизмы.

Это даёт:

- строгую математическую базу для GRA-обнулёнки,
- язык для описания субъект/инструментной двойственности,
- основу для дальнейших формализаций в GRA-экосистеме.

В минимуме:

- \(\mathcal{N} \equiv \text{хиральный подъём}\),
- хиральность = нативный язык GRA-обнулёнки.

\bigskip

\noindent
\textbf{Финальное утверждение:}

\[
\boxed{
\mathcal{N} \equiv \text{хиральный подъём}, \qquad
\text{хиральность = нативный язык GRA-обнулёнки}.
}
\]

\end{document}